\honest{Absence of Evidence Is Evidence of Absence}{Eliezer Yudkowsky}{2007}

I quote Robyn Dawes's textbook on the internment of Japanese-Americans in 1942. ``California governor Earl Warren,'' told that there had been no sabotage by Japanese-Americans, said that the lack of sabotage was ``the most ominous sign in our whole situation.'' \nb{Warren was not governor then; he was California's attorney general. The post passes the error on.}

A secret Fifth Column might delay its sabotage, I say, but no sabotage is still more likely if there is no Fifth Column at all. So the absence of sabotage counts against a Fifth Column. \nb{Warren said the sabotage was being timed like Pearl Harbor, for one coordinated day. If so, no sabotage so far was nearly as likely with a Fifth Column as without one, and it told you almost nothing.}

Then the notation, and the theorem: if seeing something would raise the probability of a hypothesis, then not seeing it must lower it.

I grant that absence of proof is not proof of absence. \nb{That is what people usually mean by the slogan the title answers, ``absence of evidence is not evidence of absence.'' The post quotes no one who uses the slogan, and the title defeats it by reading ``evidence'' in the technical sense of probability theory, not the sense its users intend.}

Then I say that absence may be ``strong evidence of absence or very weak evidence of absence,'' depending on how likely the cause was to leave signs. Gaps in the fossil record are very weak evidence, because fossils rarely form. \nb{The post does not apply this to Warren.}

I end by repeating a lesson from ``Your Strength as a Rationalist'': a model's strength lies in what it cannot explain.

\nb{The case is the imprisonment of a whole population. Apart from the quoted words of Dawes, who calls it ``a most grievous chapter,'' the post says nothing about it except which way the evidence points.}
